% Generated from locale/tr/content/lambda-calculus/lambda-definability/fixpoints.tex; do not hand-edit.


\olfileid[tr]{lam}{ldf}{fp}

\olsection{Sabit Noktalar}

Faktöriyel fonksiyonunu, aşağıdaki özelliğe sahip bir $\fn{Fac}$ terimi olarak
özyinelemeyle tanımlamak istediğimizi varsayalım:
\[
\fn{Fac} \ident \lambd[n][\fn{IsZero}\, n\, \num 1
  (\fn{Mult}\, n (\fn{Fac}(\fn{Pred} \, n)))].
\]
Yani $n=0$ ise $n$'nin faktöriyeli $1$, aksi durumda $n$ ile $n-1$'in
faktöriyelinin çarpımıdır. Elbette $\fn{Fac}$ teriminin kendisi sağ tarafta
geçtiğinden bu yolla $\fn{Fac}$'yi tanımlayamayız. Kendine gönderme içeren bu
tür özyinelemeli tanımlar lambda hesabının parçası değildir. Örneğin
\[
\fn{Mult} \ident \lambd[ab][a (\fn{Add}\, a) \num0]
\]
tanımının sağ tarafı yalnızca önceden tanımlanmış $\fn{Add}$ gibi terimleri
içerir. $\fn{Add}$'i kendi tanımlayıcı terimiyle değiştirerek her zaman
kaldırabiliriz. Böylece $\fn{Mult}$ saf bir lambda terimine dönüşür;
$\fn{Add}$ de tanımlanmış terimler içeriyorsa bu işlemi sürdürüp sonunda saf
bir lambda terimine ulaşırız.

Ancak yukarıdaki $\fn{Fac}$ tanımı gibi özyinelemeli tanımlarda bu doğru
değildir. Sağ taraftaki $\fn{Fac}$ geçişini kendi tanımıyla değiştirirsek:
\begin{align*}
  \fn{Fac} & \ident \lambd[n][\fn{IsZero}\, n\, \num{1}] \\
    & \qquad (\fn{Mult}\, n
  ((\lambd[n][\fn{IsZero} \, n \, \num 1\, (\fn{Mult}\, n\, (\fn{Fac}
    (\fn{Pred}\, n)))]) (\fn{Pred}\, n)))
\end{align*}
elde edilir ve sağ taraftaki $\fn{Fac}$ hâlâ ortadan kalkmamıştır. Bu işlemi
yinelersek tanım sürekli uzar ve hiçbir zaman saf lambda terimine ulaşmaz.
Dolayısıyla faktöriyeli, daha genel olarak özyinelemeli fonksiyonları bu
şekilde tanımlamak işe yaramaz.

Yine de özyinelemeli tanım bize bir şey söyler. $f$ faktöriyel fonksiyonunu
gösteren bir terim olsaydı
\[
\fn{Fac}' \ident \lambd[g][\lambd[n][\fn{IsZero} \, n \, \num 1\,
  (\fn{Mult} \, n\, (g (\fn{Pred} n)))]]
\]
teriminin $f$'ye uygulanması, yani $\fn{Fac}'f$, yine faktöriyel
fonksiyonunu gösterirdi. $\fn{Fac}'$yi bir fonksiyon alıp fonksiyon döndüren
bir fonksiyon olarak görürsek, $f$ faktöriyel olduğu sürece
$\fn{Fac}'f$'nin değeri yine $f$'dir. $\fn{Fac}'f\equal[\beta]f$ özelliğine
sahip bir $f$ terimine $\fn{Fac}'$nin \emph{sabit noktası} denir. Demek ki
faktöriyel $\fn{Fac}'$nin bir sabit noktasıdır.

Lambda hesabında verilen bir terimin sabit noktalarını hesaplayan terimler
vardır; bunlar $\fn{Fac}'$ gibi bir terimi faktöriyel tanımına dönüştürmekte
kullanılabilir.

\begin{defn}
  \ollabel{defn:Turing-Y}
  \emph{Y-kombinatörü} şu terimdir:
  \[
  Y \ident (\lambd[ux][x(uux)])(\lambd[ux][x(uux)]).
  \]
\end{defn}

\begin{thm}
  Her $g$ terimi için $Yg\red g(Yg)$. Dolayısıyla $Yg$ her zaman $g$'nin
  bir sabit noktasıdır.
\end{thm}

\begin{proof}
  $(\lambd[ux][x(uux)])$ terimini $U$ ile kısaltalım; böylece $Y\ident UU$.
  O zaman
  \begin{align*}
    Y g &\ident (\lambd[ux][x(uux)])U\,g \\
    &\red (\lambd[x][x(UUx)])g\\
    & \red g(UUg) \ident g(Yg).
  \end{align*}
  $g(Yg)$ ile $Yg$ aynı $g(Yg)$ terimine indirgendiğinden
  $g(Yg)\equal[\beta]Yg$ ve dolayısıyla $Yg$, $g$'nin sabit noktasıdır.
\end{proof}

$Yg$ bir redex olduğundan indirgeme elbette sonsuza kadar sürebilir:
\begin{align*}
  Y g &\red g (Y g) \\
    &\red g (g (Y g)) \\
    &\red g(g (g (Y g)))\\
    &\ldots
\end{align*}
Bu nedenle $Yg$'yi $g$'nin kendisine sonsuz kez uygulanması gibi
düşünebiliriz. $g$'yi bir kez daha uygularsak, söz gelişi, yeni bir şey
yapmış olmayız: $g$'nin $Yg$'ye sonsuz kez uygulanmasına bir $g$ daha eklemek
yine aynı sonsuz uygulamadır.

$Yg$ ile başlayan yukarıdaki $\beta$-indirgeme dizisinin sonsuz olduğuna
dikkat edin. Bu yüzden $Yg$'yi bir $N$ terimine uygularsak $(Yg)N$ de
$(g(Yg))N$, $(g(g(Yg)))N$, \dots{} gibi sonsuz sayıda farklı terime
indirgenir. Yine de \emph{başka} bir indirgeme dizisinin normal biçimde sona
ermesi mümkündür.

Faktöriyeli ele alalım. $\fn{Fac}$ terimini $Y\,\fn{Fac}'$, yani
$\fn{Fac'}$nin bir sabit noktası olarak tanımlayın. O zaman:
\begin{align*}
  \fn{Fac}\,\num 3
  &\red Y\, \fn{Fac}' \, \num 3 \\
  &\red \fn{Fac}' (Y \,\fn{Fac}') \, \num 3 \\
  &\ident (\lambd[x][\lambd[n][\fn{IsZero} \, n\, \num 1\,
      (\fn{Mult}\, n\, (x (\fn{Pred}\, n)))]]) \, \fn{Fac} \, \num 3\\
  & \red \fn{IsZero} \, \num 3\, \num 1\,
      (\fn{Mult}\, \num 3\, (\fn{Fac} (\fn{Pred}\, \num 3))) \\
  &\red  \fn{Mult}\, \num 3 \, (\fn{Fac} \, \num 2).
  \intertext{Benzer biçimde,}
  \fn{Fac}\,\num 2
  &\red  \fn{Mult}\, \num 2 \, (\fn{Fac} \, \num 1) \\
  \fn{Fac}\,\num 1
  &\red  \fn{Mult}\, \num 1 \, (\fn{Fac} \, \num 0)
  \intertext{ancak}
  \fn{Fac}\,\num 0
  &\red \fn{Fac}' (Y \,\fn{Fac}') \, \num 0 \\
  &\ident (\lambd[x][\lambd[n][\fn{IsZero} \, n\, \num 1\,
      (\fn{Mult}\, n\, (x (\fn{Pred}\, n)))]]) \, \fn{Fac} \, \num 0\\
  & \red \fn{IsZero} \, \num 0\, \num 1\,
      (\fn{Mult}\, \num 0\, (\fn{Fac} (\fn{Pred}\, \num 0)))\\
  &\red \num 1.
  \intertext{Dolayısıyla birlikte}
  \fn{Fac}\,\num 3
  &\red  \fn{Mult}\, \num 3 \,
  (\fn{Mult} \, \num 2\, (\fn{Mult} \, \num 1\, \num 1)).
\end{align*}

$\fn{Fac'}$ için geçerli olan her özyinelemeli tanım için geçerlidir.
\begin{align*}
g\,x_1\dots x_n & \equal[\beta] N
\intertext{özyinelemeli eşitliği verilsin; burada $N$, $g$ ve
$x_1$, \dots,~$x_n$ içerebilir. O zaman her zaman
$G\ident(Y\lambd[g][\lambd[x_1\dots x_n][N]])$ olacak biçimde bir $G$ vardır ve}
G\,x_1 \dots x_n & \equal[\beta] \Subst{N}{G}{g}.
\intertext{Gerçekten sabit nokta teoremi gereği}
G \ident (Y \lambd[g][\lambd[x_1\dots x_n][N]]) & \red
\lambd[g][\lambd[x_1\dots x_n][N]](Y \lambd[g][\lambd[x_1\dots x_n][N]])\\
& \ident (\lambd[g][\lambd[x_1\dots x_n][N]])\,G
\intertext{ve dolayısıyla}
G\,x_1 \dots x_n
& \red (\lambd[g][\lambd[x_1\dots x_n][N]])\,G\,x_1\dots x_n\\
& \red (\lambd[x_1\dots x_n][\Subst{N}{G}{g}])\,x_1\dots x_n\\
& \red \Subst{N}{G}{g}.
\end{align*}

\olref{defn:Turing-Y} içindeki $Y$ kombinatörü Alan Turing'e aittir. Alonzo
Church, $Y_C$ diyeceğimiz farklı bir sürüm önermiştir:
\[
Y_C \ident \lambd[g][(\lambd[x][g(xx)])(\lambd[x][g(xx)])].
\]
Church'ün kombinatörü Turing'inkinden biraz daha zayıftır: $Y_Cg$ için
$Y_Cg\equal[\beta]g(Y_Cg)$ geçerlidir, fakat $Y_Cg\bred g(Y_Cg)$ olmak
zorunda değildir. $V=\lambd[x][g(xx)]$ olsun; böylece
$Y_C\ident\lambd[g][VV]$. O zaman
\begin{align*}
  VV & \ident (\lambd[x][g(xx)])V \red g(VV)
  \text{ ve dolayısıyla}\\
  Y_C g & \ident (\lambd[g][VV])g \red VV \red g(VV), \text{ ancak ayrıca}\\
  g(Y_C g) & \ident g((\lambd[g][VV])g) \red g(VV).
\end{align*}
Başka bir deyişle $Y_Cg$ ile $g(Y_Cg)$ ortak $g(VV)$ terimine indirgenir;
dolayısıyla $Y_Cg\equal[\beta]g(Y_Cg)$. Uygulamalar için bu çoğu zaman
yeterlidir.

